% #############################################################################
% This is Chapter 5 - Research Problem
% !TEX root = ../main.tex
% #############################################################################
% Kept lean and deliberately non-duplicative: the empirical evidence for the
% claims below is presented in Chapter 1 and Chapter 4 and is referenced rather
% than restated, which is what allows this chapter to be two pages.
% #############################################################################
\fancychapter{Research Problem}
\label{chap:problem}
\label{sec:problem}

The review reported in \Cref{chap:slr} shows \ac{AI} and \acp{LLM} adopted across sectors and across every phase of the \ac{SDLC}, in a field evolving fast enough for new models and tools to appear continuously, and shows those benefits accompanied by challenges that obstruct successful implementation. The most consequential, and the most consistently reported, is the absence of established, clear methodologies shaping the proper use of these tools in professional software development~\cite{Banh:2025cf, Farhanna:2025eu, Hemmat:2025rd, Gao:2025cc}.

Several frameworks incorporating \ac{AI} into the development process have been proposed and are summarised in \Cref{tab:methodologies}; as \Cref{chap:intro} establishes, they share no common foundation regarding lifecycle integration, management practice or evaluation criteria~\cite{Ahmed:2025ai}. Organisations therefore face difficulty not only in adopting AI-supported practice but in assessing its impact on software quality, security and maintainability, no standardised benchmarks or performance metrics existing against which to do so: the field is still searching for structured and repeatable integration patterns rather than converging on them.

As \Cref{chap:intro} also establishes, established methodologies, including Waterfall, Scrum, Kanban and the Agile Manifesto, were designed for environments in which every contributor is human~\cite{Zhang:2025ea}, and fall short in AI-driven contexts in three specific ways. High-speed AI-assisted coding accelerates implementation while creating pressure upstream, in requirements clarification, and downstream, in testing, verification and maintenance. They provide no explicit structure for governing AI participation, for assessing model reliability, or for mitigating the resulting risks. And they assume all contributors are human, whereas a model is a new category of collaborator whose output carries no author and therefore no obvious owner.

The empirical case for treating this as a process problem rather than a capability problem is set out in \Cref{chap:intro} and rests on two independently reported findings: security performance does not improve as functional correctness improves, so generated code that works is not thereby safe~\cite{Veracode:2025gen}, and acceleration at the point of authorship relocates to human review rather than disappearing, without reaching delivery outcomes~\cite{Faros:2025par}. Both sources are industry-reported rather than peer-reviewed and are cited as such; their value lies in agreeing on the mechanism rather than in the precision of any single figure, neither security nor delivery improving as a by-product of faster generation.

\noindent\textbf{The core research problem addressed in this work is therefore the lack of patterned guidance, methodological rigour and standardised evaluation metrics for the systematic integration and management of \ac{AI} and \ac{LLM}-based tools across the \ac{SDLC}, while preserving human oversight, software quality and security, and organisational accountability. This gap limits the ability of organisations to realise the benefits of AI-assisted software development sustainably.}
